
\subsubsection{4.1 Research Questions}

\begin{itemize}
\item \textbf{RQ1}: Does measurable accommodation exist? (H-E1)
\item \textbf{RQ2}: Does AI adapt to human formality? (H-M2)
\item \textbf{RQ3}: Does accommodation predict engagement? (H-M3)
\item \textbf{RQ4}: Is accommodation bidirectional? (H-M1)
\end{itemize}

\subsubsection{4.2 Dataset Statistics}

\begin{table}[htbp]
\centering
\resizebox{\columnwidth}{!}{%
\begin{tabular}{ll}
\toprule
Analysis & Sample Size \\
\midrule
H-E1 (BCS) & 26,395 conversations \\
H-M1 (Lag) & 26,405 conversations \\
H-M2 (Correlation) & 111,039 turn pairs \\
H-M3 (Tercile) & 327,977 turn pairs \\
\bottomrule
\end{tabular}
}
\end{table}

\subsubsection{4.3 Evaluation Metrics}

\begin{itemize}
\item \textbf{H-E1}: BCS standard deviation (target > 0.15)
\item \textbf{H-M1}: Lag-1 correlation (target r > 0, p < 0.05)
\item \textbf{H-M2}: Pearson correlation (target |r| > 0.1, p < 0.001)
\item \textbf{H-M3}: Continuation rates by tercile
\end{itemize}

\subsubsection{4.4 Statistical Methods}

\begin{itemize}
\item Pearson and Spearman correlations for continuous relationships
\item One-sample t-tests for lag-1 correlation significance
\item Cluster bootstrap (n=2,000) for robust p-values in tercile analysis
\item Permutation tests (n=1,000) for baseline validation
\end{itemize}

